\chapter{E-Posta Guvenligi}

\section*{Giris}
E-posta halen kurumsal saldirilarin birincil giris vektorudur. Bu nedenle altyapi guvenligi, kimlik dogrulama kayitlari ve kullanici farkindaligi birlikte yonetilmelidir.

\section{Kurumsal Mesajlasma Altyapilari}
\subsection{Exchange, M365, Zimbra ve Postfix Sistemlerinde Guvenlik Mimarisi}
Mesajlasma platformlarinda MFA, rol ayrimi, SMTP relay sinirlari ve denetim loglari temel guvenlik bilesenleridir.

\section{E-Posta Dogrulama ve Itibar Yonetimi}
\subsection{SPF, DKIM ve DMARC Yapilandirmalari}
SPF, DKIM ve DMARC birlikte uygulanmadiginda spoofing/BEC riski belirgin artar.

\section{Gelismis E-Posta Tehditleri ve Savunma}
\subsection{Guvenli E-Posta Ag Gecitleri (SEG)}
SEG, anti-phishing, malware sandbox ve URL rewrite kontrollerini merkezi uygular.

\subsection{Is Sureci Ihlali (BEC) ve Oltalama Simulasyonu Araclari (Orn: GoPhish) Entegrasyonu}
BEC senaryolarinda yonetsel onay zinciri, finans akislari icin ikinci kanal dogrulama ve periyodik phishing simulasyonu kritik onemdedir.
